\chapter{Geothermal Mass Accelerators}
\label{ch:21}

The last of the principal terrestrial concepts is the one that connects the terrestrial program to the orbital one. The geothermal mass accelerator appears throughout the author's earlier notes as a ground-based system that uses the heat of the Earth as its source of energy to accelerate a payload, so that a rocket may dispense with some or all of its own propulsion in the first stage of ascent. In those notes the concept has several recurring features. It is to provide the initial thrust for vehicles that carry no engine of their own or only a small one, often through staged rings that accelerate the payload in sequence. It is to operate at low acceleration, so that it can serve passengers and fragile freight as well as hardened payloads, and the notes speak of rocketless launch and of slow freight. It is to be sited near the equator, where the rotation of the Earth contributes to the velocity of an eastward launch, and in one version on a floating structure described as an international space bridge. And it is to be integrated with the orthodromic network of Chapter 17, by retrofitting the great-circle corridors with accelerator track, and with the skyhook and expandable towers of the next part, which would receive the payload at altitude. The concept is accordingly a node in a larger architecture, and it is assessed here as such, with the assumption that the geothermal element supplies energy and the accelerator is an electromagnetic or equivalent device for converting that energy into the kinetic energy of the payload. The resource base for geothermal energy and its engineering limits are assessed in the standard interdisciplinary study~\cite{tester2006}; no source cited there addresses the accelerator, which remains the book's own proposal.

The first quantity to settle is how much energy is required, because the answer governs the whole assessment and, as it turns out, is favorable. To place a kilogram in a circular orbit at four hundred kilometers altitude requires its kinetic energy at the orbital speed of about $7.67$ km/s, which is $29.4$ MJ, and its potential energy at that altitude, about $3.8$ MJ, a total of $33.2$ MJ or $9.2$ kWh. If the payload is instead given the speed of about $8.03$ km/s that a surface-grazing transfer orbit to that altitude requires, with a small circularization burn of about $0.12$ km/s at the top, the kinetic energy is $32.2$ MJ, or $8.95$ kWh per kilogram. These figures are the energy delivered to the payload, and they exclude the losses of the accelerator and of the passage through the atmosphere. At a conversion efficiency of one half, the delivery of ten thousand tonnes a year would require about $180$ GWh a year, a mean electrical power of about $20$ MW, which a single geothermal field of modest size supplies. A field of one gigawatt, operating at ninety per cent of capacity, could in principle supply the energy for something like $880$ kilotonnes a year at perfect efficiency, which exceeds present annual launch mass to orbit by some two orders of magnitude. The energy is thus not the binding constraint, and this should be said plainly because discussions of launch often concentrate on cost per kilogram as though it were an energy cost; at the price of industrial electricity the energy for a kilogram is of the order of a dollar, and the present cost of launch is dominated by the vehicle, its propellant, and its operations.

The constraints that do bind are the track and the atmosphere. The track length required to reach a speed $v$ at a constant acceleration $a$ is $v^2/(2a)$, and for $v=8.03$ km/s it is $33$ km at $100\,g$, which is acceptable for hardened payloads and not for people, $3{,}286$ km at $1\,g$, $1{,}095$ km at $3\,g$, and $32{,}900$ km at $0.1\,g$. The last figure is of the order of the circumference of the Earth, and it explains the interest of the notes in rings and in the equator: a gentle launch implies a very long track, and the only structures that long on Earth are those that follow its curvature. Launching eastward from the equator reduces the velocity that the track must supply by the equatorial rotation speed of about $0.465$ km/s and the track length by about eleven per cent, a useful but modest gain. On the other hand the power demanded is proportional to the acceleration, because the instantaneous power of a track delivering force $ma$ at speed $v$ is $mav$, which has the mean value $\tfrac12mva$ for a launch from rest, so that a gentle launch also needs a small power plant. For a payload of one tonne, the mean power at $3\,g$ is about $118$ MW (the peak is twice as great), within the reach of a large geothermal plant and a small storage buffer, and at $0.1\,g$ the mean is about $4$ MW, whereas at $100\,g$ it is $3.9$ GW, which would have to be drawn from storage. The compatibility of geothermal baseload with low acceleration is a real feature of the concept: the form of the energy source and the form of the launch match.

The atmosphere imposes the severest constraint. A body moving at $8$ km/s at sea level meets air with a dynamic pressure $\tfrac12\rho v^2$ of nearly $40$ MPa, which is about four hundred times atmospheric pressure, and the stagnation enthalpy $\tfrac12v^2$ of $32$ MJ per kilogram would, if converted to heat in the air, correspond to temperatures of tens of thousands of kelvins before dissociation and radiation intervene. Both are conditions of reentry, run in reverse, and they would destroy any payload not protected by a shield of a mass comparable to that of reentry vehicles, only here the exposure occurs at the lowest and densest part of the atmosphere, where the loads are greatest. The dynamic pressure falls to $23.7$ MPa at five kilometers, $2.9$ MPa at twenty, and $0.13$ MPa at forty, so that high speed is bearable only at great altitude. There are two responses, both of which are present in the notes. One is to evacuate the track and extend the evacuated tube to high altitude, so that the payload leaves the tube above most of the atmosphere; the mass and strength of a tube a few tens of kilometers tall exceed anything built, and its association in the notes with expandable towers reflects this requirement. The other is to reduce the exit speed, so that the accelerator supplies only the first part of the ascent and a skyhook, a rocket, or both supply the rest. At an exit speed of $2$ km/s the dynamic pressure at sea level is $2.5$ MPa and at five kilometers $1.5$ MPa, which is about what high-speed projectiles and missiles endure.

The reduced-speed option has an attractive consequence in the mathematics of the rocket. The velocity increment that a rocket must supply to reach orbit, including losses to drag and gravity, is about $9.4$ km/s. With a specific impulse of $350$ s, corresponding to an exhaust velocity of $3.43$ km/s, the mass ratio is $e^{9.4/3.43}\approx15.5$. If the accelerator supplies $2$ km/s of the increment the mass ratio falls to $e^{7.4/3.43}\approx8.6$, and the fraction of initial mass that is not propellant rises from $6.5$ per cent to $11.6$ per cent, a factor of $1.79$. The accelerator supplies only $(2/7.67)^2\approx7$ per cent of the orbital kinetic energy, but because propellant must lift the remaining propellant, the benefit in payload fraction is out of proportion to the energy share. This is the logic of the staged arrangement: a modest accelerator, fed with inexpensive energy, makes a difference to the economics of the rocket, and a more ambitious one makes a larger difference at disproportionately higher difficulty.

The siting of the device raises the geothermal question. Geothermal resources are concentrated along tectonic boundaries and volcanic arcs, and the high-temperature fields lie in Iceland, East Africa, the Andes, Central America, Indonesia, the Philippines, and the western United States. Few are near the equator at low altitude, and those that are, such as the fields of the East African Rift and of the volcanic arcs of Ecuador and Indonesia, are inland and mountainous, which is favorable for exit altitude and unfavorable for the transport of the vehicles and for the length of a straight track. A floating equatorial installation, as in the notes, would not be located on a geothermal resource at all and would receive its energy by cable. The designation geothermal is thus best read as a statement of the energy source, which can be transmitted over distance, and not of the location of the accelerator.

A final and serious consideration is that the device is symmetric with respect to its use. An apparatus that can accelerate a mass to orbital speed can also accelerate a mass to hypersonic speed along a trajectory that does not enter orbit, and a kinetic projectile of ordinary material carries energy comparable to that of conventional explosives per unit mass at a speed of only $3$ km/s. A mass accelerator is accordingly a dual-use technology of a particularly stark kind, and its governance cannot be separated from its engineering. The institutional requirements include transparency of the facility's operations, the inspection of payloads, the registration of launches, and a treaty framework that does not exist at present. The notes describe the geothermal mass accelerators in one narrative as having been dormant and then becoming tools of sustainable growth and not weapons or instruments of control, and the statement is a good statement of the requirement, not a description of the means by which it would be met. The assessment under the standard of Chapter 16 is that the energy balance is favorable and the principal difficulties, the atmospheric exit and the length of the gentle track, are unsolved; the concept belongs to the third readiness class and must not figure in any plan on which present obligations depend.

\section*{Formalization}

Let a payload of mass $m$ be accelerated at constant acceleration $a$ along a track to a speed $v$ relative to the track, the track itself moving at an inertial speed $v_0$ along the direction of launch (the equatorial rotation, for an eastward launch).

\begin{proposition}
(i) The track length is $L=(v-v_0)^2/(2a)$ with $v$ the inertial exit speed, and the duration is $t=(v-v_0)/a$. (ii) The work done on the payload by the track is $W=\tfrac12m(v^2-v_0^2)$, and the instantaneous power is $P(s)=ma\,(v_0+as)$ at time $s$, whose time average is $\bar P=W/t=\tfrac12 m a (v+v_0)$ and whose peak is $mav$. (iii) For orbital insertion at radius $r_a$ from a surface transfer ellipse, the exit speed at the surface is $v_p=\sqrt{\mu(2/R-1/a_t)}$ with $a_t=(R+r_a)/2$, and the apogee circularization increment is $\Delta v=\sqrt{\mu/r_a}-\sqrt{\mu(2/r_a-1/a_t)}$.
\end{proposition}

\begin{proof}
(i) Kinematics in the frame of the track: $v-v_0=at$ and $L=\tfrac12 a t^2$. (ii) The inertial speed of the payload at time $s$ is $v_0+as$, the force is $ma$, and power is force times speed. Integration over $[0,t]$ gives $W=ma(v_0t+\tfrac12at^2)=\tfrac12m(v^2-v_0^2)$. The average is $W/t$, and $v_0+at=v$ gives the peak. (iii) The vis-viva equation $v^2=\mu(2/r-1/a_t)$ evaluated at the perigee radius $R$ and at the apogee radius $r_a$, with the circular speed $\sqrt{\mu/r_a}$.
\end{proof}

Numerically, with $\mu=3.986\times10^5\ \mathrm{km^3\,s^{-2}}$, $R=6371$ km, and $r_a=6771$ km, one finds $v_p=8.029$ km/s, apogee speed $7.555$ km/s, circular speed $7.673$ km/s, and $\Delta v=0.118$ km/s. The kinetic energy at $v_p$ is $32.2$ MJ kg$^{-1}$, which is $8.95$ kWh kg$^{-1}$. The track lengths at $v_0=0$ are $32{,}865$ km at $0.1\,g$, $6{,}573$ km at $0.5\,g$, $3{,}286$ km at $1\,g$, $1{,}095$ km at $3\,g$, $329$ km at $10\,g$, and $32.9$ km at $100\,g$, with durations $136$, $27$, $13.6$, $4.5$, $1.4$ minutes and $8.2$ s respectively; at the equatorial $v_0=0.465$ km/s the lengths are reduced by the factor $(1-v_0/v)^2=0.887$. The mean power $\tfrac12mav$ for $m=1000$ kg, $v=8.03$ km/s, $v_0=0$ is $3.9$ MW at $0.1\,g$, $39$ MW at $1\,g$, $118$ MW at $3\,g$, and $3.9$ GW at $100\,g$, and the peak $mav$ is twice these values.

The pressure loads are $q=\tfrac12\rho v^2$, equal to $39.5$ MPa at sea level ($\rho=1.225\ \mathrm{kg\,m^{-3}}$) for $v=8.03$ km/s, $23.7$ MPa at $5$ km ($0.736$), $2.87$ MPa at $20$ km ($0.0889$), $0.59$ MPa at $30$ km ($0.0184$), and $0.13$ MPa at $40$ km ($0.0040$).

For the benefit of partial acceleration, let the required ideal velocity increment be $\Delta v_{\mathrm{tot}}$, the accelerator supply $\Delta v_a$, and the rocket exhaust velocity $c$.

\begin{proposition}
The ideal mass ratio of the rocket stage is $\mathcal M=\exp[(\Delta v_{\mathrm{tot}}-\Delta v_a)/c]$, and relative to no accelerator the ratio of payload fractions is $\exp(\Delta v_a/c)$, which is independent of $\Delta v_{\mathrm{tot}}$. The fraction of orbital kinetic energy supplied by the accelerator is $(\Delta v_a/v_{\mathrm{orb}})^2$ when the payload starts at rest, so the payload-fraction gain exceeds the energy share whenever $\exp(\Delta v_a/c)>(\Delta v_a/v_{\mathrm{orb}})^2$.
\end{proposition}

\begin{proof}
The first statement is the rocket equation, and the ratio follows by division. The second is the ratio of $\tfrac12\Delta v_a^2$ to $\tfrac12v_{\mathrm{orb}}^2$.
\end{proof}

For $\Delta v_{\mathrm{tot}}=9.4$ km/s and $c=3.43$ km/s: $\Delta v_a=0,1,2,3$ km/s gives mass ratios $15.5$, $11.5$, $8.6$, $6.4$ and payload-fraction factors $1$, $1.34$, $1.79$, $2.40$ (ideal, ignoring structural mass), against energy shares of $0$, $1.7$, $6.8$, and $15.3$ per cent. The disproportion between energy share and payload gain is the mathematical content of the claim that a modest accelerator is worthwhile if it can be built at acceptable cost, and the same proposition shows why the claim weakens once structural mass, which scales with the stage and not with the propellant alone, is included.
